\section{Nucleo e DNA}

\subsection{Il nucleo}

All'interno del nucleo si trovano:
\begin{itemize}
  \item \textbf{cromatina}: fibre di DNA associate a proteine;
  \item uno o più \textbf{nucleoli}: strutture elettrondense di forma irregolare in cui avviene la sintesi dell'RNA ribosomiale e l'assemblaggio dei ribosomi;
  \item \textbf{nucleoplasma}: sostanza fluida in cui sono disciolti i soluti del nucleo;
  \item \textbf{matrice nucleare}: rete fibrillare contenente proteine.
\end{itemize}
Il nucleo è delimitato da una doppia membrana, l'\textbf{involucro nucleare}, che presenta dei \textbf{pori nucleari}.

% ---------------------------------------------------------------------------
\subsection{Involucro nucleare}

Costituito da 2 membrane cellulari parallele, distanti 10--50\,nm. È una barriera che impedisce il passaggio di ioni, soluti e macromolecole tra nucleo e citoplasma; è attraversato da \textbf{pori}, costituiti da proteine, nei punti in cui le due membrane si fondono.

\begin{itemize}
  \item \textbf{membrana esterna}: in continuità con il reticolo endoplasmatico rugoso (RER);
  \item \textbf{membrana interna}: legata a una sottile rete fibrillare (\textbf{lamina nucleare}) mediante proteine integrali di membrana.
\end{itemize}

% ---------------------------------------------------------------------------
\subsection{Pori nucleari}

Permettono il passaggio di RNA e proteine in entrambe le direzioni. Sono costituiti da un complesso apparato a forma di canestro, il \textbf{complesso del poro nucleare} (CPN), che sporge sia verso il citoplasma che verso il nucleoplasma.

\begin{itemize}
  \item simmetria ottagonale, dovuta alla ripetizione per 8 volte di specifiche strutture;
  \item formato da circa 30 proteine differenti, ciascuna ripetuta 8 volte.
\end{itemize}

\rubrica{Trasporto attraverso il poro nucleare}
Le proteine che presentano un segnale di localizzazione nucleare (NLS) si legano a un recettore eterodimerico (\textbf{importina} $\alpha/\beta$), formando un complesso che si associa a un filamento citoplasmatico. Il complesso recettore-carico si muove attraverso il poro nucleare ed entra nel nucleoplasma, dove interagisce con Ran-GTP e si dissocia. La subunità $\beta$ dell'importina, associata a Ran-GTP, viene trasportata indietro nel citoplasma, dove il Ran-GTP viene idrolizzato e quindi riconvertito a Ran-GTP; l'importina $\alpha$ viene trasportata indietro nel citoplasma da un'\textbf{esportina}.

% ---------------------------------------------------------------------------
\subsection{Cromatina e cromosomi}

\rubrica{Cromatina}
Il DNA nel nucleoplasma è associato a istoni e ad alcune proteine a formare la \textbf{cromatina}. Nelle cellule non in divisione si presenta sotto forma di granuli dispersi e può assumere due forme:
\begin{itemize}
  \item \textbf{eterocromatina}: forma condensata, in grossi ammassi irregolari;
  \item \textbf{eucromatina}: forma dispersa, sotto forma di piccoli granuli o filamenti.
\end{itemize}
Nelle cellule in divisione la cromatina si condensa in piccoli bastoncini, detti \textbf{cromosomi}, capaci di autoreplicazione e di mantenere le proprie caratteristiche morfologiche attraverso successive divisioni cellulari.

\rubrica{Nucleosoma}
Gli \textbf{istoni} sono carichi positivamente ed è quindi possibile la loro associazione con il DNA, carico negativamente, a formare il \textbf{nucleosoma}: primo livello di organizzazione della cromatina.
\begin{itemize}
  \item ciascun nucleosoma è costituito da 8 istoni (ottamero istonico: 2 H2A + 2 H2B + 2 H3 + 2 H4) e da 146 paia di basi di DNA che si avvolgono attorno ad essi $\rightarrow$ diametro di circa 10\,nm;
  \item dall'ottamero sporgono le \textbf{code istoniche};
  \item tra un nucleosoma e l'altro è presente il \textbf{DNA di giunzione} (o \textit{DNA linker}), di lunghezza variabile (da 8 a 114 coppie di nucleotidi), associato a un ulteriore istone (\textbf{H1});
  \item la successione di nucleosomi collegati dal DNA linker forma una struttura a ``collana di perle'', del diametro di circa 11\,nm;
  \item i nucleosomi adiacenti possono compattarsi tra loro per formare una fibra del diametro di circa 30\,nm;
  \item ulteriori superavvolgimenti di questa struttura, in associazione a proteine non istoniche, portano alla formazione dei cromosomi visibili durante la metafase della divisione mitotica.
\end{itemize}

\rubrica{Struttura del cromosoma}
Ogni cromosoma contiene una molecola di DNA a doppio filamento, singola e continua.
\begin{itemize}
  \item ciascuna delle 2 estremità contiene un tratto di brevi sequenze ripetute associate a proteine specializzate a formare un ``cappuccio'' $\rightarrow$ il \textbf{telomero};
  \item il cromosoma duplicato è formato da 2 unità parallele identiche $\rightarrow$ i \textbf{cromatidi};
  \item lungo il cromosoma è presente una zona ristretta (costrizione primaria), detta \textbf{centromero};
  \item associate al centromero sono presenti delle proteine specifiche che fungono da siti di attacco per i microtubuli $\rightarrow$ i \textbf{cinetocori}.
\end{itemize}

\rubrica{Classificazione dei cromosomi}
In base alla posizione del centromero:
\begin{center}
\begin{forest} classalbero
[Cromosomi
  [metacentrico]
  [submetacentrico]
  [acrocentrico]
  [telocentrico]
]
\end{forest}
\end{center}

% ---------------------------------------------------------------------------
\subsection{Struttura del DNA}

Il DNA è costituito da una doppia elica formata da due filamenti antiparalleli, ciascuno composto da unità ripetute di zucchero, fosfato e base azotata. Le basi dei due filamenti sono appaiate al centro dell'elica mediante legami idrogeno; lungo l'elica si distinguono un solco maggiore e un solco minore.

% ---------------------------------------------------------------------------
\subsection{Replicazione del DNA}

\rubrica{Principio: appaiamento di basi complementari}
\begin{itemize}
  \item appaiamento tra basi complementari: A--T, G--C;
  \item i due filamenti si svolgono e si separano;
  \item ciascun filamento funge da stampo per l'aggiunta di basi delle nuove eliche;
  \item il risultato sono 2 nuove doppie eliche di DNA identiche a quella parentale, ciascuna formata da un filamento ``vecchio'' e uno ``nuovo'' $\rightarrow$ \textbf{replicazione semiconservativa}.
\end{itemize}

\rubrica{Filamento guida e filamento in ritardo}
Le DNA polimerasi si muovono lungo lo stampo solo in direzione $3' \rightarrow 5'$. Di conseguenza, i due filamenti di nuova sintesi crescono in direzioni opposte:
\begin{itemize}
  \item \textbf{filamento guida}: assemblato in modo continuo, nella stessa direzione della forcella di replicazione;
  \item \textbf{filamento in ritardo}: assemblato in segmenti (frammenti di Okazaki), successivamente uniti tra loro.
\end{itemize}

\rubrica{Fasi della sintesi del filamento in ritardo}
\begin{enumerate}
  \item sintesi del primer a RNA da parte della \textbf{primasi};
  \item allungamento da parte della \textbf{DNA polimerasi III};
  \item rimozione del primer e riempimento delle interruzioni da parte della \textbf{DNA polimerasi I};
  \item fusione dei filamenti da parte della \textbf{DNA ligasi}.
\end{enumerate}

\rubrica{Proteine ed enzimi coinvolti alla forcella di replicazione}
\begin{itemize}
  \item \textbf{elicasi}: srotola il DNA alla forcella di replicazione;
  \item \textbf{girasi} (topoisomerasi II): taglia la molecola duplex e ne rilassa la tensione meccanica dovuta allo srotolamento, introducendo tagli a doppio filamento a monte della forcella e facendo ruotare le estremità tagliate attorno all'asse centrale, scaricando così lo stress torsionale;
  \item \textbf{proteine che legano il DNA a singolo filamento} (SSB): legano e stabilizzano il DNA a singolo filamento alla forcella di replicazione;
  \item \textbf{complesso dell'RNA primasi} (primosoma): l'RNA primer inizia la sintesi di un nuovo filamento;
  \item \textbf{complesso della DNA polimerasi};
  \item \textbf{DNA ligasi}: unisce i frammenti di Okazaki sul filamento lento.
\end{itemize}

\rubrica{Completamento del filamento in ritardo (rimozione dei primer a RNA)}
\begin{itemize}
  \item la \textbf{DNA polimerasi delta} allunga il filamento a monte fino a raggiungere il primer a RNA del frammento precedente del filamento lento;
  \item una proteina che lega il DNA a singolo filamento (\textbf{RPA}) sposta l'RNA e una piccola porzione di DNA;
  \item la proteina RPA attira una specifica endonucleasi che taglia il frammento spostato;
  \item il primer a RNA e il DNA tagliati vengono degradati da esonucleasi;
  \item la DNA polimerasi delta prosegue la sintesi e rimpiazza l'RNA primer; l'ultimo legame è formato dalla DNA ligasi.
\end{itemize}

% ---------------------------------------------------------------------------
\subsection{Errori nella replicazione}

La riparazione di un accoppiamento di basi errato (\textit{mismatch}) consiste nell'escissione di un piccolo frammento di DNA contenente la base appaiata in modo erroneo, seguita dalla sintesi del frammento corretto.
\begin{itemize}
  \item uno dei due filamenti viene scelto per essere riparato;
  \item un segmento del filamento da riparare viene escisso ed eliminato;
  \item il gap rimanente viene colmato durante la riparazione, usando l'altro filamento come stampo;
  \item i due possibili tipi di riparazione (a seconda del filamento scelto) portano a differenti sequenze di DNA.
\end{itemize}
